\section*{Hoofdstuk \ref{ch:b1:active-margins} --- Actieve randen}

\begin{solution}{exo:b1:active-margins:1}
(a) $2\pi G\,\Delta\rho\,h = 2\pi\times\num{6.674e-11}\times1700\times3000 =
\qty{2.14e-3}{m.s^{-2}}$: de zwaartekracht zou met ongeveer
\qty{214}{\milli\gal} lager zijn. (b) Voor \qty{3.5}{km},
$2\pi\times\num{6.674e-11}\times1700\times3500 = \qty{2.50e-3}{m.s^{-2}}$,
\qty{250}{\milli\gal}: de grootte van de waargenomen laagte, dus de trog wordt ver van isostatisch evenwicht gehouden. (c) De formule van de plaatzijde gaat uit van gesteente dat aan alle kanten zonder grens doorloopt; een trog die slechts tientallen kilometers breed is, verwijdert veel minder massa rond het meetpunt, dus het water alleen geeft minder dan \qty{250}{\milli\gal}. Het overige tekort ligt dieper: de hele plaat is onder de trog naar beneden gebogen, en haar \omterm{def:b1:layered-earth:layers}{korst}, lichter dan de \omterm{def:b1:layered-earth:layers}{mantel}, neemt aan haar basis de plaats van de \omterm{def:b1:layered-earth:layers}{mantel} in.
\end{solution}

\begin{solution}{exo:b1:active-margins:2}
(a) $\qty{82}{mm.yr^{-1}} = \qty{82}{km.Myr^{-1}}$, dus \qty{820}{km} in
\qty{10}{Myr}. (b) $820\sin\ang{45} = \qty{580}{km}$. (c) De weg naar
\qty{600}{km} is $600/\sin\ang{45} = \qty{849}{km}$, afgelegd in
$849/82 = \qty{10.3}{Myr}$.
\end{solution}

\begin{solution}{exo:b1:active-margins:3}
$g = 110/\tan\delta$: \qty{191}{km} voor \ang{30}, \qty{110}{km} voor \ang{45},
\qty{64}{km} voor \ang{60}. De boog boven de flauw hellende plaatzijde ligt het verst van zijn trog.
\end{solution}

\begin{solution}{exo:b1:active-margins:4}
(a) Convergentie: $\qty{2.7e6}{m^2.yr^{-1}}/\qty{5.1e7}{m} =
\qty{0.053}{m.yr^{-1}}$, ongeveer \qty{53}{mm.yr^{-1}}. Spreiding:
$\qty{3.0e6}{m^2.yr^{-1}}/\qty{6.7e7}{m} = \qty{0.045}{m.yr^{-1}}$, ongeveer
\qty{45}{mm.yr^{-1}}. Over een kortere totale lengte verbruiken de subductiegrenzen de oceaanbodem sneller dan de ruggen haar maken, en toch verbruiken zij in totaal minder.
(b) Het overige \qty{0.3}{km^2} of zo wordt opgenomen langs de andere convergente grenzen van het model, zoals continentale botsingen en convergente zones binnen de oceanen, waar de platen worden verkort in plaats van gesubduceerd; en de ruggen zijn niet de enige divergerende grenzen: ook continentale riften en sommige transforms openen een beetje.
\end{solution}

\begin{solution}{exo:b1:active-margins:5}
(a) $\bar x = 350$, $\bar z = 290$; de afwijkingen van $x$ zijn $-150$, $-50$,
$50$, $150$ en die van $z$ zijn $-140$, $-40$, $40$, $140$, dus
$S_{xx} = 2(150^2 + 50^2) = \num{50000}$ en
$S_{xz} = 2(150\times140 + 50\times40) = \num{46000}$. Dan $s = 0.92$,
$a = 290 - 0.92\times350 = \qty{-32}{km}$, en de helling is
$\arctan 0.92 = \ang{42.6}$. (b) $z = 0$ bij $x = 32/0.92 = \qty{34.8}{km}$
vanaf de trog. (c) Door de eindpunten is de helling $280/300 = 0.933$, een
helling van \ang{43.0}: dichtbij, maar het gebruikt slechts twee van de vier aardbevingen.
\end{solution}

\begin{solution}{exo:b1:active-margins:6}
(a) $50/\tan\ang{25} + 60/\tan\ang{45} = 107.2 + 60.0 = \qty{167.2}{km}$.
(b) Een plaatzijde van \ang{45} overal geeft \qty{110}{km}: het flauwe ondiepe segment verbreedt de kloof met \qty{57}{km}. (c) Het diepe segment is
$z = 50 + (x - 107.2)$; het bereikt $z = 0$ bij $x = \qty{57.2}{km}$, landwaarts van de trog, zoals de Koerilese lijn doet bij \qty{53}{km}: een flauw ondiep segment verklaart de Koerilese verschuiving. (d) $107.2 + 60/\tan\ang{60} = 107.2 + 34.6 =
\qty{141.9}{km}$.
\end{solution}

\begin{solution}{exo:b1:active-margins:7}
(a) De straal is de helft van de helling, en $r = R\,\theta$ met $\theta$ in radialen:
\ang{22.5} en $6371\times0.3927 = \qty{2502}{km}$ voor \ang{45}; \ang{15} en
\qty{1668}{km} voor \ang{30}; \ang{35} en \qty{3892}{km} voor \ang{70}.
(b) De bogen boven flauw hellende plaatzijden. (c) De voorspelling is
$44.9/2 = \ang{22.45}$; de passing, \ang{21.9}, is
$0.55/22.45 = \qty{2.4}{\percent}$ kleiner.
\end{solution}

\begin{solution}{exo:b1:active-margins:8}
(a) $q/k$: $0.035/2.5 = \qty{0.014}{K.m^{-1}} = \qty{14}{K.km^{-1}}$ onder
de \omterm{def:b1:active-margins:active-margin}{voorboog}, $0.080/2.5 = \qty{32}{K.km^{-1}}$ onder het \omterm{def:b1:active-margins:active-margin}{achterboogbekken}.
(b) $10 + 14\times30 = \qty{430}{\celsius}$ en
$10 + 32\times30 = \qty{970}{\celsius}$. (c) Een deel van de oppervlaktewarmtestroom wordt geproduceerd door radioactiviteit binnen de \omterm{def:b1:layered-earth:layers}{korst} zelf, dus komt er minder warmte op de diepere niveaus aan en neemt de gradiënt met de diepte af: een lineaire extrapolatie van de oppervlaktegradiënt overschat de temperatuur op diepte.
\end{solution}

\begin{solution}{exo:b1:active-margins:9}
(a) \qty{32}{mm.yr^{-1}} is \qty{32}{km.Myr^{-1}}; de halve snelheid is
\qty{16}{mm.yr^{-1}}. (b) $100/32 = \qty{3.1}{Myr}$. (c) De snelheid hoeft niet constant te zijn geweest; een deel van de breedte kan komen van het uitrekken en riften van de boog voordat de spreiding begon;
en de as kan zijn verschoven of asymmetrisch zijn gespreid.
\end{solution}

\begin{solution}{exo:b1:active-margins:10}
(a) Het pad is $600/\sin\ang{45} = \qty{849}{km}$, afgelegd met
\qty{82}{km.Myr^{-1}} in \qty{10.3}{Myr}, dat is \qty{3.27e14}{s}.
(b) $\sqrt{\num{e-6}\times\num{3.27e14}} = \qty{1.8e4}{m}$, ongeveer
\qty{18}{km}. (c) $(\qty{5e4}{m})^2/\qty{e-6}{m^2.s^{-1}} = \qty{2.5e15}{s}$,
ongeveer \qty{79}{Myr}. (d) Warmte heeft alleen de buitenste \qty{18}{km} aan elke kant van de plaatzijde bereikt: haar kern bereikt \qty{600}{km} bijna even koud als toen zij de trog verliet, kouder en sterker dan de \omterm{def:b1:layered-earth:layers}{mantel} eromheen, en daarom kan zij nog breken in aardbevingen op die diepte.
\end{solution}

\begin{solution}{exo:b1:active-margins:11}
(a) Op afstand $L$ achter de teen staat het oppervlak $L\tan\alpha$ boven het niveau van de teen en ligt de basis $L\tan\beta$ eronder: de achterkant van het prisma is $L(\tan\alpha + \tan\beta)$ hoog, en de driehoek heeft de oppervlakte
$\frac12 L\times L(\tan\alpha + \tan\beta)$. (b) $\tan\ang{4} + \tan\ang{6} =
0.175$. Het prisma ontvangt $\qty{1}{km}\times\qty{40}{km.Myr^{-1}} =
\qty{40}{km^2.Myr^{-1}}$ doorsnede, dus $A = 40t$ en
$L = \sqrt{2A/0.175}$: \qty{21.4}{km} na \qty{1}{Myr}, \qty{42.8}{km}
na \qty{4}{Myr}, \qty{67.6}{km} na \qty{10}{Myr}. $L$ groeit als
$\sqrt t$: vier keer zo lang geeft tweemaal de breedte. (c) $A$ wordt gehalveerd en $L$ gedeeld door $\sqrt2$: \qty{15.1}{km}, \qty{30.2}{km} en \qty{47.8}{km}.
\end{solution}

\begin{solution}{exo:b1:active-margins:12}
(a) $\partial F/\partial a = -2\sum(z_i - a - s x_i) = 0$ geeft
$n\bar z = na + sn\bar x$, dat is $\bar z = a + s\bar x$. (b) Elke $z_i$ en
$\bar z$ groeien met \qty{10}{km}, dus de afwijkingen $z_i - \bar z$ en $S_{xz}$
zijn ongewijzigd: de helling is dezelfde, en $a = \bar z - s\bar x$ groeit met
\qty{10}{km}. (c) De afwijkingen $x_i - \bar x$ zijn ongewijzigd, dus de helling is
dezelfde; $\bar x$ daalt met 20, dus $a$ neemt toe met $20s$. (d) Geen van beide:
de helling hangt alleen van de richtingscoëfficiënt af. Fouten die alle punten gemeen hebben, verschuiven de lijn zonder haar te draaien.
\end{solution}

\begin{solution}{pb:b1:active-margins:1}
\textbf{1.} De acht punten liggen dicht bij een rechte lijn die onder ongeveer
\ang{45} helt, van \qty{96}{km} diep op \qty{155}{km} van de trog tot
\qty{573}{km} diep op \qty{609}{km}.

\textbf{2.} $\sum x_i = 3152$, $\sum z_i = 2751$, $\sum x_i^2 = \num{1442910}$,
$\sum x_i z_i = \num{1290548}$; $\bar x = \qty{394.0}{km}$,
$\bar z = \qty{343.875}{km}$.

\textbf{3.} $S_{xx} = 1442910 - 8\times394^2 = \num{201022}$ en
$S_{xz} = 1290548 - 8\times394\times343.875 = \num{206654}$. Dan
$s = 206654/201022 = 1.028$ en $a = 343.875 - 1.02802\times394 =
\qty{-61.2}{km}$: $z = -61.2 + 1.028\,x$.

\textbf{4.} $\arctan1.028 = \ang{45.8}$.

\textbf{5.} Op $x = 61.2/1.028 = \qty{59.5}{km}$ landwaarts van de trog. Een plaatzijde die van de trog af met deze helling begon, zou daar het oppervlak bereiken; zij bereikt het aan de trog zelf, dus helt zij nabij de trog flauwer en wordt zij daaronder steiler.

\textbf{6.} De afwijkingen zijn $-2.1$, $2.0$, $-7.3$, $5.3$, $7.1$, $2.3$,
$-15.0$ en $8.1$ \unit{km}. Aardbeving 7 ligt \qty{15}{km} boven de lijn, aardbeving 8
\qty{8}{km} eronder; de andere liggen binnen ongeveer \qty{7}{km} ervan, een breuk van de dikte van de plaatzijde.

\textbf{7.} $\ang{320.9} - \ang{300} = \ang{20.9}$. Loodrecht op de
trog: $87.4\cos\ang{20.9} = \qty{81.6}{mm.yr^{-1}}$; evenwijdig:
$87.4\sin\ang{20.9} = \qty{31.2}{mm.yr^{-1}}$. De convergentie is schuin.

\textbf{8.} \qty{81.6}{mm.yr^{-1}} is \qty{81.6}{km.Myr^{-1}}: \qty{816}{km} in
\qty{10}{Myr}.

\textbf{9.} Het pad is $632/\sin\ang{45} = \qty{894}{km}$, afgelegd in
$894/81.6 = \qty{11.0}{Myr}$.

\textbf{10.} $122 + 11 = \qty{133}{Myr}$.

\textbf{11.} $t = \qty{10.95}{Myr} = \qty{3.46e14}{s}$ en
$\sqrt{\num{e-6}\times\num{3.46e14}} = \qty{1.86e4}{m}$, ongeveer \qty{19}{km}
van elke wand: veel minder dan de \qty{50}{km} tot het midden van een plaatzijde van
\qty{100}{km} dik, waarvan de kern nog koud is.

\textbf{12.} Hier is $q/k$: \qty{17.2}{K.km^{-1}} onder de oude oceaanbodem,
\qty{14.0}{K.km^{-1}} onder de \omterm{def:b1:active-margins:active-margin}{voorboog} en \qty{32.0}{K.km^{-1}} onder het
\omterm{def:b1:active-margins:active-margin}{achterboogbekken}.

\textbf{13.} Onder de \omterm{def:b1:active-margins:active-margin}{voorboog} ligt de koude plaatzijde slechts enkele tientallen kilometers diep; zij neemt warmte op uit de rots erboven, dus bereikt er minder warmte het oppervlak.

\textbf{14.} $h = \qty{2.60e-3}{m.s^{-2}}/(2\pi\times\num{6.674e-11}\times1700)
= \qty{3647}{m}$, ongeveer \qty{3.6}{km}: even veel als, of iets meer dan, de
hele diepte van de trog onder de bodem eromheen, ongeveer \qty{3.5}{km}. De
trog mist massa die door niets wordt gecompenseerd.

\textbf{15.} Het dichtheidsverschil is $3300 - 1000 = \qty{2300}{kg.m^{-3}}$,
en $2\pi\times\num{6.674e-11}\times2300\times3500 = \qty{3.38e-3}{m.s^{-2}}$,
\qty{338}{\milli\gal}. De formule van de plaatzijde telt de ontbrekende massa tot op elke afstand; de echte trog is slechts over een band van tientallen kilometers breed gebogen, dus haar tekort is kleiner, ongeveer \qty{260}{\milli\gal}, tussen de twee schattingen van de plaatzijde in.

\textbf{16.} De sterkte van de gebogen plaat, die de trog omlaag houdt en de buitenrug omhoog zoals een elastische plaat dat zou doen, en de trek van de dichte plaatzijde die in de \omterm{def:b1:layered-earth:layers}{mantel} zinkt.

\textbf{17.} $-61.2 + 1.028\times164 = \qty{107.4}{km}$.

\textbf{18.} Ja: \qty{107}{km} ligt tussen 72 en \qty{173}{km}, en slechts
ongeveer \qty{2}{km} dieper dan het gemiddelde, \qty{105}{km}.

\textbf{19.} $x = (100 + 61.2)/1.028 = \qty{156.8}{km}$.

\textbf{20.} $\ang{45.8}/2 = \ang{22.9}$, en
$6371\times22.9\times\pi/180 = \qty{2546}{km}$.

\textbf{21.} De passing, \ang{21.9}, is \ang{1.0} kleiner, ongeveer
\qty{4}{\percent}: de gebogen bol beschrijft de kromming van de Koerilese boog goed.

\textbf{22.} $x = (105 + 61.2)/1.028 = 166.2/1.028 = \qty{161.7}{km}$.

\textbf{23.} $164 - 161.7 = \qty{2.3}{km}$ dichter bij de trog dan het waargenomen
front, $2.3/164 = \qty{1.4}{\percent}$: de Koerilese vulkanen staan boven een plaatzijde die heel licht dieper is dan het gemiddelde, zoals vraag 18 vond.

\textbf{24.} Tot op de kilometer nauwkeurig voorspelt de lijn van de Koerilese plaatzijde een
\omterm{def:b1:active-margins:volcanic-arc}{boog--trog-afstand} van \textbf{\qty[text-series-to-math=true]{162}{km}} voor
vulkanen die boven de plaatzijde op \qty{105}{km} staan (waargenomen: \qty{164}{km}).
\end{solution}
